% ================================================================
% v5 SKELETON — restart after v4 review verdict: 4.8/10, major flaw
% diagnosed as structural, not a bug to patch.
%
% DIAGNOSIS CARRIED OVER FROM v4 REVIEW (do not re-litigate without
% re-reading it): a fixed-time, L-uniformly-bounded Julia-channel
% construction is, by its own uniform-boundedness hypothesis,
% *incapable* of seeing the one thing that defines NHSE: an
% extensive, boundary-condition-dependent effect. Its resource was
% positive for gauge-trivial Hermitian phase chains and insensitive
% to OBC vs PBC. That is not a gap to shrink with more lemmas; it is
% the wrong object class.
%
% DESIGN REQUIREMENT LIST for v5 (a "universal QRT of NHSE" must,
% at minimum):
%   (R1) be EXTENSIVE in L under OBC for genuinely skin-effect models,
%   (R2) VANISH (or stay O(1)) under PBC for the same Hamiltonian,
%   (R3) vanish under any operation that is "obviously free" physically
%        (local reciprocal recoupling, global unitary relabeling),
%   (R4) reduce to a computable closed form on the Hatano-Nelson family,
%   (R5) connect to the GBZ / non-Bloch band theory for general
%        translation-invariant models, not just the 2x2/HN toy cases,
%   (R6) be numerically checkable independently of any fixed-time
%        approximation.
% v4 failed (R1) and (R2) outright and never attempted (R5). This
% skeleton is organized so every section states which requirement it
% targets.
%
% HOW TO USE THIS SKELETON: every \begin{todo}...\end{todo} block is
% unproved / unimplemented. Sections without a todo block, or with
% "SALVAGED FROM v4" / "NEW, PROVED BELOW" tags, are usable as-is.
% Do not delete the diagnosis block above when revising.
% ================================================================
\documentclass[reprint,amsmath,amssymb,aps,nofootinbib]{revtex4-2}

\usepackage{graphicx}
\usepackage{bm}
\usepackage{mathtools}
\usepackage{mathrsfs}
\usepackage{amsthm}
\usepackage{bbm}
\usepackage{hyperref}
\usepackage{microtype}
\usepackage{booktabs}
\usepackage{xcolor}

\newtheorem{definition}{Definition}
\newtheorem{theorem}{Theorem}
\newtheorem{proposition}{Proposition}
\newtheorem{lemma}{Lemma}
\newtheorem{corollary}{Corollary}
\newtheorem{remark}{Remark}
\newtheorem{assumption}{Assumption}
\newtheorem{conjecture}{Conjecture}

% Visible TODO environment -- renders in red so it cannot be missed
% or accidentally left in a submitted draft.
\newenvironment{todo}
  {\begin{quote}\color{red}\bfseries TODO (unproved/unimplemented):\normalfont\itshape}
  {\end{quote}}

\newcommand{\Tr}{\operatorname{Tr}}
\newcommand{\id}{\mathbbm{1}}
\newcommand{\GBZ}{\mathrm{GBZ}}
\newcommand{\NHSE}{\mathrm{NHSE}}
\newcommand{\diag}{\operatorname{diag}}
\newcommand{\norm}[1]{\left\lVert #1\right\rVert}
\newcommand{\spr}{\operatorname{spr}}
\newcommand{\cond}{\operatorname{cond}}
\newcommand{\Rres}{\mathcal R}
\newcommand{\Free}{\mathfrak F}
\newcommand{\GaugeC}{\mathfrak G}

\begin{document}
\preprint{APS/123-QED}

\title{%
% TODO: title must not promise "NHSE" unless (R1)+(R2) are actually
% delivered with a proof, not just the HN worked example.
A Resource Theory of Non-Hermitian Skin Effect from Extensive
Non-Normality: [working title]
}

\author{Ann Author}
\affiliation{Authors' institution and/or address}
\date{\today}

\begin{abstract}
% TODO: do not write the final abstract until Secs.~\ref{sec:gbz}
% and \ref{sec:lower} are either proved or explicitly demoted to
% conjectures in the abstract's own language. The v4 abstract's
% failure mode was overclaiming a biconditional that only held on
% one family; do not repeat it.
[Abstract placeholder. Skeleton claim to eventually support:] We
define a resource theory of the non-Hermitian skin effect in which
the resource is the logarithmic condition number of the minimal-cost
similarity transformation bringing a non-Hermitian propagator to a
reciprocal (transpose-symmetric) form. Unlike a fixed-time channel
distance, this quantity is extensive in system size under open
boundary conditions for skin-effect models and is exactly zero,
uniformly in size, under periodic boundary conditions for the same
Hamiltonian family, reproducing the boundary-condition sensitivity
that defines NHSE. [Remainder depends on
Secs.~\ref{sec:gbz}--\ref{sec:lower}.]
\end{abstract}

\maketitle

% ================================================================
\section{Introduction: what a universal NHSE resource theory must do}
\label{sec:intro}

Quantum/operator resource theories fix an object class, a free
subset, free transformations, and a preorder
\cite{ChitambarGour,LiuYuan,ChiribellaSupermap}. A resource theory
of NHSE additionally inherits a non-negotiable empirical constraint:
NHSE is, by definition, a phenomenon of \emph{open} boundary
conditions that vanishes under periodic boundary conditions for the
identical bulk Hamiltonian, and whose characteristic scale (the skin
depth) is a bulk quantity while its manifestation (macroscopic
eigenvector accumulation) is extensive in system size
\cite{HatanoNelson,YaoWang,OkumaEtAl2020}. Any resource measure
$\Rres$ built from an $L$-uniformly bounded generator evaluated at a
fixed time is a Lipschitz function of the Hamiltonian on the
relevant compact set and therefore cannot reproduce an
OBC/PBC discontinuity or extensive scaling; this is not a proof
gap, it is a dimensional-analysis obstruction. A prior attempt by
the present authors built exactly such an object (fixed-time Julia
channel dilation, diamond-norm asymmetry) and, on review, was shown
to (i) assign positive resource to a gauge-trivial Hermitian phase
chain with no skin effect, and (ii) assign the same resource to a
Hamiltonian under OBC and PBC, i.e., it failed both (R1) and (R2)
above. That construction is retained below only as a local,
sub-extensive, experimentally-motivated diagnostic
(Appendix~\ref{app:v4}), never as the primary monotone.

The resource we use instead is well precedented outside the QRT
literature: NHSE is characterized by extensive \emph{non-normality}
of the OBC Hamiltonian, i.e., by the fact that the transformation
diagonalizing $H_L$ (or bringing it to any reciprocal/normal form)
becomes exponentially ill-conditioned in $L$
\cite{OkumaEtAl2020,TrefethenEmbree}. We package this as a resource
theory: the free objects are reciprocal (transpose-symmetric)
operators, the free operations are similarity transformations of
$L$-independent cost, and the monotone is the log-condition-number
of the cheapest transformation to a free object.

\begin{todo}
Literature check: confirm citation for "extensive non-normality
$\Leftrightarrow$ NHSE" as an existing folklore/technical result
(candidates: Okuma--Sato reviews, Trefethen--Embree pseudospectra,
Longhi's non-Hermitian imaging papers) and state precisely how the
present resource-theoretic packaging differs from that prior work,
so the novelty claim in the abstract is accurate.
\end{todo}

% ================================================================
\section{Object class, free set, and the conditioning resource}
\label{sec:setup}

\begin{definition}[Object class]
For each $L$, the object class is
$\mathsf{Obj}_L=\{H\in\mathbb C^{L\times L}\}$, no boundedness or
normalization imposed a priori beyond finiteness. Physical families
are indexed sequences $\{H_L\}_{L\ge1}$, typically with $H_L$ the
open-boundary truncation of a fixed local (finite-range) coupling
rule.
\end{definition}

\begin{definition}[Free set: reciprocal operators]
\begin{equation}
\Free_L=\{H\in\mathsf{Obj}_L: H^T=H\}.
\end{equation}
\end{definition}
This is the same reciprocity notion used in v4 (transpose symmetry
in the fixed physical basis), retained deliberately for continuity:
what changes in v5 is not the notion of ``free,'' but the object
class and the cost functional. Note $\Free_L$ is not restricted to
Hermitian matrices; diagonal matrices with complex entries are
reciprocal, which is essential below (Sec.~\ref{sec:hn-pbc}).

\begin{definition}[Cheap gauge group]
For $C\ge1$,
\begin{equation}
\GaugeC_{L,C}=\{S\in GL(L,\mathbb C):\cond(S)\le C\},
\qquad
\cond(S)=\norm S\,\norm{S^{-1}},
\end{equation}
acting by $H\mapsto SHS^{-1}$. We call $\GaugeC_{L,C}$ \emph{free}
for fixed, $L$-independent $C$: it represents any basis change,
local relabeling, or bounded-strength local recoupling whose cost
does not grow with system size.
\end{definition}

\begin{definition}[Reciprocity-conditioning resource]
\label{def:Rres}
\begin{equation}
\boxed{
\Rres(H)=\inf\{\log\cond(S): S\in GL(L,\mathbb C),\ SHS^{-1}\in\Free_L\}
\in[0,\infty].}
\end{equation}
\end{definition}

\begin{proposition}[Approximate monotonicity under free gauge]
\label{prop:monotone}
% NEW, PROVED BELOW -- this replaces v4's monotonicity theorem with
% one that is nontrivial precisely because the object class is now
% extensive.
For every $S_0\in\GaugeC_{L,C}$,
\begin{equation}
\boxed{\Rres(S_0HS_0^{-1})\le\Rres(H)+2\log C.}
\end{equation}
\end{proposition}

\begin{proof}
Let $S$ be any competitor for $H$ with $SHS^{-1}\in\Free_L$. Then
$(SS_0^{-1})(S_0HS_0^{-1})(SS_0^{-1})^{-1}=SHS^{-1}\in\Free_L$, and
submultiplicativity of the condition number gives
$\cond(SS_0^{-1})\le\cond(S)\cond(S_0^{-1})=\cond(S)\cond(S_0)\le
C\cond(S)$. Take the infimum over $S$.
\end{proof}

\begin{remark}
This is an \emph{additive-constant} monotonicity, not the strict
monotonicity of v4. That is the correct notion for an extensive
resource: it says free operations of bounded cost can change
$\Rres$ by at most an $L$-independent amount, so the leading,
extensive-in-$L$ part of $\Rres$ is a genuine invariant of the free
orbit. This mirrors how resource theories of asymptotic/extensive
quantities (e.g. entanglement cost) are normally stated.
\end{remark}

\begin{proposition}[Diagonalizable upper bound]
\label{prop:diag-bound}
% NEW, PROVED BELOW.
If $H=V\Lambda V^{-1}$ with $\Lambda$ diagonal, then
$\Rres(H)\le\log\cond(V)$.
\end{proposition}
\begin{proof}
Diagonal matrices are transpose-symmetric, so $\Lambda\in\Free_L$;
take $S=V^{-1}$.
\end{proof}

\begin{remark}
Proposition~\ref{prop:diag-bound} identifies $\Rres$ as bounded by
the standard eigenvector-matrix condition number used in
pseudospectral analyses of non-normal operators. This is the
precise sense in which the new resource is ``non-normality,''
resource-theoretically packaged.
\end{remark}

% ================================================================
\section{Worked example: Hatano--Nelson, OBC vs PBC}
\label{sec:hn-pbc}

\subsection{Open boundary conditions: extensive resource}
% SALVAGED FROM v4 (Sec.~VI there): the similarity gauge $S_g$ was
% already derived and is reused verbatim, but now as the actual
% object of study rather than as a step toward an eigenvector
% formula for a different quantity.

Let $H_{g,L}=m\id+Je^gT_++Je^{-g}T_-$ on the open chain, $T_-=T_+^T$
the nearest-neighbor shift. With $S_g=\diag(e^{-gn})_{n=1}^L$,
\begin{equation}
S_gH_{g,L}S_g^{-1}=H_{0,L},
\end{equation}
which is real symmetric, hence reciprocal: $H_{0,L}\in\Free_L$.

\begin{proposition}[OBC upper bound]
\label{prop:hn-obc}
% NEW, PROVED BELOW.
\begin{equation}
\boxed{\Rres(H_{g,L})\le|g|(L-1).}
\end{equation}
\end{proposition}
\begin{proof}
$\cond(S_g)=\max_n e^{-gn}/\min_n e^{-gn}=e^{|g|(L-1)}$; apply the
definition with $S=S_g$.
\end{proof}

This is extensive in $L$, giving (R1) for the first time in this
line of work, with an explicit, non-asymptotic constant.

\subsection{Periodic boundary conditions: exact triviality}
\label{subsec:pbc}
% NEW, PROVED BELOW. This is the key structural fix: it is the
% first result in this research program that actually distinguishes
% OBC from PBC, which v4 could not do even in principle.

Under PBC, $H_{g,L}^{\rm PBC}$ is circulant and is diagonalized by
the (unitary) discrete Fourier matrix $F$, $\cond(F)=1$.

\begin{proposition}[PBC exact vanishing]
\label{prop:hn-pbc}
\begin{equation}
\boxed{\Rres(H_{g,L}^{\rm PBC})=0\quad\text{for every }L\text{ and every }g.}
\end{equation}
\end{proposition}
\begin{proof}
$F H_{g,L}^{\rm PBC}F^{-1}=\diag(\varepsilon_g(k))_k$ is diagonal,
hence reciprocal, and $\cond(F)=1$ (unitary), so
$\Rres\le\log1=0$; also $\Rres\ge0$ always.
\end{proof}

\begin{corollary}[First genuine OBC/PBC separation in this program]
\label{cor:separation}
For $g\ne0$,
\begin{equation}
\Rres(H_{g,L}^{\rm OBC})-\Rres(H_{g,L}^{\rm PBC})
\ \text{grows at least like}\ |g|(L-1),
\end{equation}
while the identical bulk coupling rule is used in both cases. This
delivers (R1) and (R2) simultaneously and is the property v4 never
achieved.
\end{corollary}

\begin{todo}
The lower bound $\Rres(H_{g,L}^{\rm OBC})\ge c\cdot|g|L$ for some
$c>0$ (matching Prop.~\ref{prop:hn-obc} up to a constant) is NOT yet
proved; Prop.~\ref{prop:hn-obc} only gives an upper bound, i.e. it
shows the resource is \emph{at most} extensive via one particular
gauge, not that no cheaper gauge exists. Likely strategy: use the
known asymptotic eigenvector localization
$\psi_{j,g}(n)\propto e^{-gn}\sin(\pi jn/(L+1))$ (already in v4,
Sec.~VI) together with a Bauer--Fike / pseudospectral argument
lower-bounding $\cond(V)$ for \emph{any} diagonalizing $V$, then
extend from ``diagonalizing'' competitors to \emph{all} competitors
$S$ with $SHS^{-1}$ merely reciprocal (not necessarily diagonal).
The second step is the harder one and is completely open.
\end{todo}

% ================================================================
\section{General translation-invariant chains and the GBZ}
\label{sec:gbz}
% TARGETS (R5). Nothing here is proved; this section is the actual
% scientific content still missing and is the difference between a
% one-model paper and a genuine "universal" theory.

\begin{todo}
State and prove (or falsify) a conjecture of the form
\begin{equation}
\lim_{L\to\infty}\frac{\Rres(H_L^{\rm OBC})}{L}
=2\int (\text{something built from }\ln|\beta(k)|\text{ on the GBZ}),
\end{equation}
i.e. that the extensive rate of the conditioning resource equals
(twice) the GBZ radius / a $k$-averaged non-Bloch decay rate, for
general translation-invariant nearest-neighbor (or finite-range)
non-Hermitian chains $H(\beta)=\sum_j H_j\beta^j$, not just the
scalar Hatano--Nelson case. The natural candidate gauge to try first
is the generalized exponential similarity
$S_\rho=\diag(\rho^{-n})$ for $\rho=e^{\bar\mu}$ with $\bar\mu$ set
by the GBZ radius, generalizing $S_g$ above; the open step is
proving this is asymptotically optimal (matching upper and lower
bounds) for general multi-band, multi-root GBZ geometry, not only
the single-root scalar case where it is exact by construction.
Likely needs the non-Bloch band theory machinery of
\cite{OkumaEtAl2020,YaoWang} and a pseudospectral/Bauer--Fike bound
as in the lower-bound TODO above. This is the load-bearing result
that would justify the word "universal" in the title; without it,
the paper is a (better) worked example, not a general theory.
\end{todo}

% ================================================================
\section{Free superchannels and a genuine conversion theory}
\label{sec:free-ops}
% TARGETS (R3), and revisits the one thing v4 could state cleanly
% (a conversion preorder) but never made physically compelling.

\begin{todo}
v4's free-superchannel class was axiomatic (transpose-covariance of
an abstract superchannel) with no laboratory cost model attached.
Here the free operations (Sec.~\ref{sec:setup}) are already more
physical -- bounded-cost basis changes / local recoupling -- but
the following are still needed:
\begin{enumerate}
\item A precise physical dictionary: which laboratory operations on
  a coupled non-Hermitian lattice (gain/loss retuning, local
  hopping edits) correspond to conjugation by some $S\in\GaugeC_{L,C}$
  with $C$ independent of $L$? Right now $\GaugeC_{L,C}$ is defined
  algebraically; it needs an operational reading before referees
  will accept it as "free," matching the same criticism v4 received
  about its axiomatic covariance class.
\item A conversion theorem: necessary and/or sufficient conditions
  for $H\to H'$ (both OBC families) under $\GaugeC_{\cdot,C}$ for
  some fixed $C$, at least on the HN gauge orbit, analogous to the
  $|g|\ge|g'|$ result in v4 but for the new resource. Likely
  tractable: on the HN orbit, $\Rres(H_{g,L})\approx|g|L$
  (pending the lower bound TODO above), so conversion should reduce
  to comparing $|g|$ and $|g'|$ directly, i.e. redo v4's
  Sec.~V/Cor.~7 argument in the new resource, which should be
  easier here since the new $\Rres$ is exactly additive under
  stacking chains, unlike the old bounded diamond-norm quantity.
\end{enumerate}
\end{todo}

% ================================================================
\section{Numerical program}
\label{sec:numerics}

\begin{todo}
No script yet implements $\Rres$. Needed, in order of priority:
\begin{enumerate}
\item \textbf{OBC-vs-PBC scaling check.} For the HN chain, compute
  $\cond(V_L)$ (eigenvector matrix, via \texttt{numpy.linalg.eig}
  plus SVD-based condition number) for $H_{g,L}^{\rm OBC}$ across a
  range of $L$ and confirm $\log\cond(V_L)/L\to$ const $>0$ as
  predicted by Prop.~\ref{prop:hn-obc}, versus
  $\cond(V_L)=1$ (up to floating point) for $H_{g,L}^{\rm PBC}$, as
  predicted by Prop.~\ref{prop:hn-pbc}. This directly checks
  Cor.~\ref{cor:separation} and is the single most important new
  numerical result the paper needs -- it did not exist in any v4
  form. Reuse \texttt{hn\_chain()} from
  \texttt{nhse\_locality\_verification\_v4.py} as the OBC generator
  and add a periodic-hopping variant for PBC.
\item \textbf{Full-$\Rres$ vs.\ diagonalization-only upper bound.}
  $\Rres$ as defined optimizes over \emph{all} $S$ with
  $SHS^{-1}$ merely reciprocal, not only diagonal; the diagonalizing
  $S=V^{-1}$ is only a feasible point. A numerical lower bound on
  $\Rres$ independent of the diagonalization upper bound (e.g. via
  an SDP relaxation or a pseudospectral argument) would give
  numerical evidence for the open lower-bound TODO in
  Sec.~\ref{sec:hn-pbc}, but no such solver exists yet in this
  project.
\item \textbf{Multi-band GBZ test.} Once Sec.~\ref{sec:gbz} states a
  precise conjecture, build a second test Hamiltonian with a
  non-scalar (matrix-valued) GBZ (e.g.\ an SSH-type non-Hermitian
  chain) and check the conjectured rate numerically.
\end{enumerate}
\end{todo}

% ================================================================
\section{Conclusion and roadmap}
\label{sec:conclusion}

The reciprocity-conditioning resource $\Rres$
(Def.~\ref{def:Rres}) is, by Props.~\ref{prop:hn-obc}--\ref{prop:hn-pbc}
and Cor.~\ref{cor:separation}, the first object in this research
program that is simultaneously extensive under OBC and exactly zero
under PBC on a genuine NHSE example -- the two properties the v4
referee report identified as structurally unreachable for a
fixed-time bounded-channel construction. What remains, in order of
how load-bearing each piece is for the title's claim of
universality:
\begin{enumerate}
\item the OBC lower bound (Sec.~\ref{sec:hn-pbc}, TODO) -- without
  it, Prop.~\ref{prop:hn-obc} is only an upper bound and the
  ``extensive'' claim is unproved, only exhibited;
\item the general GBZ theorem (Sec.~\ref{sec:gbz}, TODO) -- without
  it, this is a Hatano--Nelson paper with a better resource
  definition, not a universal theory, and the title must say so;
\item an operational free-operation dictionary and a conversion
  theorem (Sec.~\ref{sec:free-ops}, TODO);
\item the numerical program (Sec.~\ref{sec:numerics}, TODO), item 1
  of which is cheap and should be done first since it is the
  paper's single most persuasive figure.
\end{enumerate}
Do not add v4-style scope-limiting remarks as a substitute for doing
1--3; the lesson from v4 was that hedging cannot repair a wrong
object class, and the corollary is that hedging also should not be
used to avoid finishing a right one.

% ================================================================
\appendix
\section{Salvaged v4 material: a local diagnostic, not the monotone}
\label{app:v4}

% SALVAGED FROM v4. The Julia-channel construction and its exact
% short-time coefficient (v4 Theorem, Sec. VII there) are
% mathematically correct and reduce to a genuinely local,
% experimentally-motivated quantity once stripped of the claim that
% it is a QRT monotone for NHSE. Demote, don't delete: reintroduce
% $\mathcal N_{H,t}$, the exact short-time slope
% $S_{\kappa_0}(H)=\tfrac12\spr(G_{\kappa_0}(H))$, and the
% measurable state-overlap witness $d_m(t)$/$\Gamma_{\rm ps}(H)$ as
% an Appendix presenting a bulk, boundary-condition-\emph{insensitive}
% companion diagnostic (explicitly note it is the same under OBC and
% PBC, in contrast with $\Rres$), useful because it is directly
% measurable without extracting the full eigenvector-conditioning
% quantity. Do not restate it as a resource monotone for NHSE; it is
% not one, per the v4 review.

\begin{todo}
Port in condensed form (not full re-derivation) from
\texttt{hnse\_qrt\_revised\_v4.tex}: the Julia completion
(Prop.~1--Cor.~1 there), the exact short-time coefficient theorem,
and the measurable witness proposition, all clearly labeled as
Appendix material illustrating a local companion diagnostic. Cut
everything else from v4 (generalized robustness, the two-site
rigidity theorem, the buffered-channel locality section) -- the
peer review found these tangential to the paper's actual claims and
they should not be reintroduced here.
\end{todo}

% ================================================================
\begin{acknowledgments}
\end{acknowledgments}

% ================================================================
\begin{thebibliography}{99}
\bibitem{ChitambarGour}
E.~Chitambar and G.~Gour,
\textit{Quantum resource theories},
Rev.~Mod.~Phys.~\textbf{91}, 025001 (2019).
\bibitem{LiuYuan}
Y.~Liu and X.~Yuan,
\textit{Operational resource theory of quantum channels},
Phys.~Rev.~Research~\textbf{2}, 012035(R) (2020).
\bibitem{ChiribellaSupermap}
G.~Chiribella, G.~M.~D'Ariano, and P.~Perinotti,
\textit{Transforming quantum operations: quantum supermaps},
Eur.~Phys.~Lett.~\textbf{83}, 30004 (2008).
\bibitem{HatanoNelson}
N.~Hatano and D.~R.~Nelson,
\textit{Localization transitions in non-Hermitian quantum mechanics},
Phys.~Rev.~Lett.~\textbf{77}, 570 (1996).
\bibitem{YaoWang}
S.~Yao and Z.~Wang,
\textit{Edge states and topological invariants of non-Hermitian systems},
Phys.~Rev.~Lett.~\textbf{121}, 086803 (2018).
\bibitem{OkumaEtAl2020}
N.~Okuma, K.~Kawabata, T.~Shiozaki, and M.~Sato,
\textit{Topological origin of non-Hermitian skin effects},
Phys.~Rev.~Lett.~\textbf{124}, 086801 (2020).
\bibitem{TrefethenEmbree}
L.~N.~Trefethen and M.~Embree,
\textit{Spectra and Pseudospectra: The Behavior of Nonnormal
Matrices and Operators} (Princeton University Press, 2005).
% TODO: this is a book citation added on general knowledge of the
% pseudospectra literature; verify edition/details before submission.
\end{thebibliography}
\end{document}
